Fix an integer \(p\geq0\), numbers \(\nu\geq2\) and \(L\geq4\), and a law \(P\in\mathcal P_{12}(p,\nu,L)\).  Temporarily write \(h=h_n\) and \(B=B_n\).  For \(t\in\{0,1\}\) and \(x\in\mathcal B_P\), let \(\beta^*_{t,P,x}(h)\) solve the original, unwinsorized population normal equations, and define
\[
 D_h=\sup_{t\in\{0,1\},\,x\in\mathcal B_P}
 \lVert\widehat\Psi_{t,x}(h)-\Psi_{t,P,x}(h)\rVert,
 \qquad
 \mathcal G_h=\{D_h\leq(2L)^{-1}\}.
\]
All expectations of possibly nonmeasurable suprema below are outer expectations.

On \(\mathcal G_h\), clause \(\mathrm{(viii)}\) of \(\mathrm{def:cty\text{-}a1\text{-}a2\text{-}class}\) and Weyl's eigenvalue inequality give, uniformly in \((t,x)\),
\[
 \lambda_{\min}\{\widehat\Psi_{t,x}(h)\}
 \geq \lambda_{\min}\{\Psi_{t,P,x}(h)\}-D_h
 \geq L^{-1}-(2L)^{-1}=(2L)^{-1}.
\]
Thus the stabilization branch does not activate and
\(\lVert\widehat\Psi_{t,x}(h)^{-1}\rVert\leq2L\).  The empirical normal equations and the definition of \(\widehat O^B_{t,x}(h)\) then imply
\[
 \widehat\beta^B_{t,x}(h)-\beta^*_{t,P,x}(h)
 =\widehat\Psi_{t,x}(h)^{-1}\widehat O^B_{t,x}(h),
\]
and hence
\[
 \sup_{t,x}\lVert\widehat\beta^B_{t,x}(h)-\beta^*_{t,P,x}(h)\rVert
 \leq 2L\sup_{t,x}\lVert\widehat O^B_{t,x}(h)\rVert
 \quad\text{on }\mathcal G_h.
\]

Subtracting the unwinsorized population normal equations from the definition of the winsorized score yields
\[
 \mathbb E_P\widehat O^B_{t,x}(h)
 ={1\over h^2}\mathbb E_P\!\left[
 \mathbf1\{D^{\pm}_{P,x}\in I_t\}K_\square(D^{\pm}_{P,x}/h)
 r_p(D^{\pm}_{P,x}/h)\{\psi_B(Y)-Y\}
 \right].
\]
On the observed arm, clause \(\mathrm{(i)}\) gives \(Y=Y(t)\).  Each coordinate of the polynomial basis is bounded by one on the kernel support.  The conditional bound in \(\mathrm{lem:cty\text{-}winsorization\text{-}bias}\), the density bound \(f_P\leq L\), and containment of the kernel support in the radius-\(h\) disk therefore give
\[
 \begin{split}
 \sup_{t,x}\lVert\mathbb E_P\widehat O^B_{t,x}(h)\rVert
 &\leq {CB^{-3}\over h^2}
 \sup_{x\in\mathcal B_P}P\{\lVert X-x\rVert_2\leq h\}\\
 &\leq {CB^{-3}\over h^2}\,L\pi h^2
 \leq CB^{-3}.
 \end{split}
\]
The centered bound in \(\mathrm{lem:cty\text{-}winsorized\text{-}score\text{-}maximal\text{-}bound}\) and the triangle inequality for outer expectation now show
\[
 \mathbb E_P^*\sup_{t,x}\lVert\widehat O^B_{t,x}(h)\rVert
 \leq C\left\{B^{-3}+\sqrt{{\log(B/h)\over nh^2}}
 +{B\log(B/h)\over nh^2}\right\}.
\]

By \(\mathrm{lem:cty\text{-}uniform\text{-}first\text{-}order\text{-}bias}\),
\[
 \sup_{x\in\mathcal B_P}\left|
 e_1^\top\{\beta^*_{1,P,x}(h)-\beta^*_{0,P,x}(h)\}-\tau_P(x)
 \right|\leq Ch.
\]
The identification of the target used here is the signed-distance identity in \(\mathrm{lem:cty\text{-}distance\text{-}identification}\).  Clause \(\mathrm{(iii)}\) gives \(|\mu_{t,P}(x)|\leq L\), and therefore \(|\tau_P(x)|\leq2L\).  Projection onto \([-2L,2L]\) is one-Lipschitz and fixes \(\tau_P(x)\).  On \(\mathcal G_h^c\), both the projected estimator and the target belong to this interval, so their absolute difference is at most \(4L\).  Combining these observations pointwise gives
\[
 \sup_{x\in\mathcal B_P}
 |\widehat\tau^{12,\mathrm{LP}}_{n,h}(x)-\tau_P(x)|
 \leq Ch+C\sup_{t,x}\lVert\widehat O^B_{t,x}(h)\rVert
 +4L\mathbf1\{D_h>(2L)^{-1}\}.
\]
Since \(\mathbf1\{D_h>(2L)^{-1}\}\leq2LD_h\), the Gram bound in \(\mathrm{lem:cty\text{-}expected\text{-}local\text{-}polynomial\text{-}maximal\text{-}bounds}\) implies
\[
 \mathbb E_P^*\mathbf1\{D_h>(2L)^{-1}\}
 \leq C\sqrt{{\log(1/h)\over nh^2}}
 \leq C\sqrt{{\log(B/h)\over nh^2}},
\]
where the last inequality uses \(B\geq1\).  Monotonicity and subadditivity of outer expectation consequently yield, uniformly over \(P\in\mathcal P_{12}(p,\nu,L)\),
\[
 \mathbb E_P^*\!\left[
 \sup_{x\in\mathcal B_P}
 |\widehat\tau^{12,\mathrm{LP}}_{n,h}(x)-\tau_P(x)|
 \right]
 \leq C\left\{h+B^{-3}+\sqrt{{\log(B/h)\over nh^2}}
 +{B\log(B/h)\over nh^2}\right\}.
\]

Set \(h=a_n=(\log n/n)^{1/4}\) and \(B=a_n^{-1/3}\).  Then \(h\downarrow0\), and the regime required by the established maximal bounds holds because
\[
 {n^{(1+\nu)/(2+\nu)}h^2\over\log(1/h)}
 \asymp {n^{\nu/\{2(2+\nu)\}}\over\sqrt{\log n}}
 \longrightarrow\infty.
\]
Moreover,
\[
 \log(B/h)={4\over3}\log(1/h)
 ={1\over3}(\log n-\log\log n),
\]
so
\[
 h=a_n,\qquad B^{-3}=a_n,
 \qquad \sqrt{{\log(B/h)\over nh^2}}=O(a_n),
\]
while
\[
 {B\log(B/h)\over nh^2}
 =O\!\left(n^{-5/12}(\log n)^{5/12}\right)=o(a_n).
\]
All constants depend only on \((p,\nu,L)\) and are uniform in \(P\) and \(n\).  Taking the supremum over \(P\in\mathcal P_{12}(p,\nu,L)\) proves the asserted bound for all sufficiently large \(n\).  Finally, \(\mathrm{def:cty\text{-}stabilized\text{-}local\text{-}polynomial\text{-}estimator}\) supplies this rule as one explicit, law-independent member of \(\mathcal T_n^{12,\pm}\).